% Kenngrößen von Verteilungen – bank/kenngroessen-von-verteilungen/ (zusammenbau v0.4, Heft)
\einheitenkopf[t6]{Kenngrößen von Verteilungen}
\setcounter{aufgabe}{26}

% Anlauf (ohne Prüfkennung)
\begin{aufgabe}{Erwartungswert berechnen und deuten – Anlauf}
\begin{teile}
\teil Ein Würfelspiel kostet 1\,€ Einsatz. $X$ ist der Betrag in Euro, den der Spieler nach dem Wurf erhält. Was beschreibt $X$? \\ \kreuz{Auszahlung} \\ \kreuz{Gewinn}
\teil Die Tabelle zeigt die Verteilung der Zufallsgröße $X$; eine Wahrscheinlichkeit fehlt. Ergänze sie und berechne den Erwartungswert $E(X)$.

\sachtabelle{lccc}{$x$ & 0 & 1 & 2}{$P(X = x)$ & 0{,}5 & 0{,}3 & \leerzelle}
fehlend: \leerfeld , $E(X) = $ \leerfeld
\teil Bei einem Losverkauf bringen 2\,\% der Lose 50\,€, 10\,\% der Lose 5\,€, die übrigen nichts. $X$ ist die Auszahlung für ein Los in Euro. Wie groß ist $E(X)$? $E(X) = $ \leerfeld[€]
\end{teile}
\end{aufgabe}

\begin{aufgabe}{Erwartungswert berechnen und deuten – Prüfungsaufgaben}
\begin{teile}
\teil Ein Spiel kostet 3\,€ Einsatz. Mit der Wahrscheinlichkeit 0{,}25 werden 10\,€ ausgezahlt, sonst nichts. Zeige, dass das Spiel für den Anbieter auf lange Sicht gewinnbringend ist, und gib seinen mittleren Gewinn je Spiel an \hfill (Abitur 2018 GK). \leerfeld[€] je Spiel
\teil Ein Spiel kostet 4\,€ Einsatz. Ausgezahlt werden 0\,€ mit der Wahrscheinlichkeit 0{,}5, 5\,€ mit 0{,}3 und 12\,€ mit 0{,}2. Jana sagt: „Der Erwartungswert der Auszahlung ist positiv, also gewinne ich auf lange Sicht.“ Stimmt das \hfill (Abitur 2023 GK)?
\teil In einer Urne liegen drei rote und sieben blaue Kugeln. Ein Spieler zieht zweimal mit Zurücklegen. Bei zwei roten Kugeln erhält er 9\,€, bei genau einer roten Kugel 4\,€, sonst nichts. Der Term $9 \cdot 0{,}3^2 + 4 \cdot 2 \cdot 0{,}3 \cdot 0{,}7 + 0 \cdot 0{,}7^2$ gibt den Erwartungswert der Auszahlung in Euro an. Erläutere die Bedeutung des zweiten Summanden im Sachzusammenhang \hfill (Abitur 2023 GK).
\teil Ein Glücksrad zeigt mit der Wahrscheinlichkeit $\frac{1}{4}$ Grün. Es wird zweimal gedreht. Bei zweimal Grün gibt es 6\,€, bei genau einmal Grün 1\,€, sonst nichts. Der Term $6 \cdot \left(\frac{1}{4}\right)^2 + 1 \cdot 2 \cdot \frac{1}{4} \cdot \frac{3}{4}$ gibt den Erwartungswert der Auszahlung in Euro an. Erläutere den zweiten Summanden und erkläre, warum die 2 darin keine Auszahlung ist \hfill (Abitur 2023 GK).
\end{teile}
\end{aufgabe}

% Anlauf (ohne Prüfkennung)
\begin{aufgabe}{Prozentuale Abweichung einer Anzahl vom Erwartungswert berechnen – Anlauf}
\begin{teile}
\teil Ein Würfel wird 300-mal geworfen, dabei fallen 42 Sechsen. Um wie viel Prozent weicht die Zahl der Sechsen vom Erwartungswert ab? \leerfeld \%
\end{teile}
\end{aufgabe}

\begin{aufgabe}{Prozentuale Abweichung einer Anzahl vom Erwartungswert berechnen – Prüfungsaufgaben}
\begin{teile}
\teil In einer Gärtnerei keimen erfahrungsgemäß 6\,\% der Samen nicht. Von 150 Samen keimen 12 nicht. Um wie viel Prozent weicht diese Anzahl vom Erwartungswert ab? Runde auf eine Stelle \hfill (Abitur 2022 GK). \leerfeld \%
\end{teile}
\end{aufgabe}

% Anlauf (ohne Prüfkennung)
\begin{aufgabe}{Rückwärts – Anlauf}
\begin{teile}
\teil Ein Würfelspiel zahlt bei einer Sechs 4\,€ aus, sonst nichts. Gefragt ist, wie viel im Mittel je Spiel ausgezahlt wird. Vorwärts (Erwartungswert berechnen) oder rückwärts (Unbekannte aus einer Erwartungswertbedingung)? \\ \kreuz{vorwärts} \\ \kreuz{rückwärts}
\teil Ein Spiel kostet 2\,€ Einsatz. Mit der Wahrscheinlichkeit 0{,}25 wird der Betrag $a$ (in Euro) ausgezahlt, sonst nichts. Für welches $a$ ist das Spiel fair? $a = $ \leerfeld[€]
\teil Ein Würfel trägt auf vier Seiten die Zahl 1 und auf zwei Seiten die Zahl $b$. $X$ ist die geworfene Zahl, und es gilt $E(X) = 3$. Wie groß ist $b$? $b = $ \leerfeld
\end{teile}
\end{aufgabe}

\begin{aufgabe}{Rückwärts – Prüfungsaufgaben}
\begin{teile}
\teil In einer Urne liegen sechs weiße und $x$ schwarze Kugeln. Zwei Kugeln werden mit einem Griff gezogen; gewonnen wird, wenn beide weiß sind. Wie viele schwarze Kugeln müssen in der Urne liegen, damit die Gewinnwahrscheinlichkeit $\frac{1}{3}$ beträgt \hfill (Abitur 2018 GK)? $x = $ \leerfeld
\teil Ein Spieler legt 3\,€ Einsatz aus und zieht zweimal mit Zurücklegen eine Kugel aus einem Beutel mit goldenen und silbernen Kugeln. Nach jedem Zug wird der ausliegende Betrag bei Gold verdreifacht, bei Silber gedrittelt; am Ende erhält er den ausliegenden Betrag. $p$ ist der Anteil der goldenen Kugeln. Für welches Verhältnis von goldenen zu silbernen Kugeln haben Spieler und Spielleiter die gleiche Gewinnerwartung \hfill (Abitur 2023 GK)? $p = $ \leerfeld , Gold : Silber $= $ \leerfeld
\end{teile}
\end{aufgabe}

% Anlauf (ohne Prüfkennung)
\begin{aufgabe}{Erwartungswertgleichung für einen Glücksradparameter aus den Spielregeln herleiten – Anlauf}
\begin{teile}
\teil Ein Glücksrad hat Felder mit den Zahlen 2 und 1; die 2 erscheint mit der Wahrscheinlichkeit $q$. Es wird zweimal gedreht; die Punktzahl ist das Produkt der beiden Zahlen. Im Mittel soll sie 2{,}25 betragen. Zeige, dass $q$ die Gleichung $4q^2 + 8q - 5 = 0$ erfüllt, und bestimme $q$. $q = $ \leerfeld
\end{teile}
\end{aufgabe}

\begin{aufgabe}{Erwartungswertgleichung für einen Glücksradparameter aus den Spielregeln herleiten – Prüfungsaufgaben}
\begin{teile}
\teil Ein Glücksrad hat Felder mit den Zahlen 4 und 1; die 4 erscheint mit der Wahrscheinlichkeit $q$. Es wird zweimal gedreht; der Rabatt in Prozent ist das Produkt der beiden Zahlen. Auf lange Sicht soll der Rabatt im Mittel 4\,\% betragen. Zeige, dass $q$ die Gleichung $3q^2 + 2q - 1 = 0$ erfüllt, und bestimme $q$ \hfill (Abitur 2024 GK). $q = $ \leerfeld
\end{teile}
\end{aufgabe}

% Anlauf (ohne Prüfkennung)
\begin{aufgabe}{Zwei Wahrscheinlichkeiten einer Verteilung aus dem Erwartungswert und der Summe 1 bestimmen – Anlauf}
\begin{teile}
\teil Die Zufallsgröße $X$ nimmt die Werte 1, 2 und 4 an, mit $P(X = 4) = 0{,}2$ und $E(X) = 2{,}2$. Bestimme $P(X = 1)$ und $P(X = 2)$. $P(X = 1) = $ \leerfeld , $P(X = 2) = $ \leerfeld
\end{teile}
\end{aufgabe}

\begin{aufgabe}{Zwei Wahrscheinlichkeiten einer Verteilung aus dem Erwartungswert und der Summe 1 bestimmen – Prüfungsaufgaben}
\begin{teile}
\teil Eine App vergibt beim täglichen Öffnen 5, 10 oder 40 Münzen; 40 Münzen gibt es weiter mit der Wahrscheinlichkeit 0{,}05. Die Wahrscheinlichkeiten für 5 und 10 Münzen werden so geändert, dass in 30 Tagen im Mittel 270 Münzen anfallen. Bestimme die beiden neuen Wahrscheinlichkeiten \hfill (Abitur 2021 GK). $P(5) = $ \leerfeld , $P(10) = $ \leerfeld
\end{teile}
\end{aufgabe}

% Anlauf (ohne Prüfkennung)
\begin{aufgabe}{Streuung – Anlauf}
\begin{teile}
\teil Ein Würfel wird 30-mal geworfen; $X$ ist die Anzahl der Sechsen. Welche Formel liefert die Standardabweichung von $X$? \\ \kreuz{Binomialformel} \\ \kreuz{allgemeine Formel}
\teil $X$ ist binomialverteilt mit $n = 100$ und dem Erwartungswert $\mu = 50$. Bestimme $p$ und die Standardabweichung $\sigma$. $p = $ \leerfeld , $\sigma = $ \leerfeld
\teil $X$ ist binomialverteilt, ihre Verteilung ist symmetrisch, und der Erwartungswert ist 7. Bestimme $p$ und $n$. $p = $ \leerfeld , $n = $ \leerfeld
\end{teile}
\end{aufgabe}

\begin{aufgabe}{Streuung – Prüfungsaufgaben}
\begin{teile}
\teil $X$ ist binomialverteilt mit dem Erwartungswert 12 und der Varianz 3. Bestimme $n$ und $p$ \hfill (Abitur 2017 GK). $n = $ \leerfeld , $p = $ \leerfeld
\teil $X$ ist binomialverteilt mit dem Erwartungswert $\mu = 12$ und der Standardabweichung $\sigma = 3$. Bestimme $p$ und $n$ und gib einen Term für $P(X = 13)$ an \hfill (Abitur 2026 GK). $p = $ \leerfeld , $n = $ \leerfeld
\teil Begründe, dass es bei 18 Versuchen keine Trefferwahrscheinlichkeit gibt, für die die Standardabweichung der Trefferzahl 2{,}5 ist \hfill (Abitur 2017 GK).
\teil Ein Graph zeigt $\sigma$ in Abhängigkeit von $p$ für binomialverteilte Zufallsgrößen mit $n = 1\,200$; die Achsen tragen Kästchen, aber keine Zahlen. Der Bogen beginnt im Ursprung und trifft nach zwanzig Kästchen wieder die $p$-Achse; über dem fünften Kästchen der $p$-Achse liegt er genau fünf Kästchen hoch. Gib an, welcher Wert zu einem Kästchen jeder Achse gehört, und erläutere dein Vorgehen \hfill (Abitur 2024 GK). $p$-Achse: \leerfeld je Kästchen, $\sigma$-Achse: \leerfeld je Kästchen
\teil Ein Graph zeigt $\sigma$ in Abhängigkeit von $p$ für binomialverteilte Zufallsgrößen mit $n = 3\,600$; die Achsen tragen Kästchen, aber keine Zahlen. Der Bogen beginnt im Ursprung und trifft nach zehn Kästchen wieder die $p$-Achse; über dem ersten Kästchen der $p$-Achse liegt er genau vier Kästchen hoch. Gib an, welcher Wert zu einem Kästchen jeder Achse gehört, und erläutere, warum man nicht den höchsten Punkt zum Skalieren nimmt \hfill (Abitur 2024 GK). $p$-Achse: \leerfeld je Kästchen, $\sigma$-Achse: \leerfeld je Kästchen
\end{teile}
\end{aufgabe}

% Anlauf (ohne Prüfkennung)
\begin{aufgabe}{Stichprobenumfang aus einer vorgegebenen Standardabweichung der Binomialverteilung berechnen – Anlauf}
\begin{teile}
\teil Eine Münze wird $n$-mal geworfen; die Anzahl von Zahl soll die Standardabweichung 5 haben. Wie groß ist $n$? $n = $ \leerfeld
\end{teile}
\end{aufgabe}

\begin{aufgabe}{Stichprobenumfang aus einer vorgegebenen Standardabweichung der Binomialverteilung berechnen – Prüfungsaufgaben}
\begin{teile}
\teil Die Trefferzahl einer Bernoulli-Kette mit der Trefferwahrscheinlichkeit 10\,\% soll die Standardabweichung 6 haben. Wie viele Versuche sind nötig \hfill (Abitur 2017 GK)? $n = $ \leerfeld
\end{teile}
\end{aufgabe}

% Anlauf (ohne Prüfkennung)
\begin{aufgabe}{Diagramm – Anlauf}
\begin{teile}
\teil Eine Münze wird zehnmal geworfen; das Diagramm zeigt die Verteilung der Anzahl von Zahl. Was ist der Erwartungswert dieser Anzahl? \\ \kreuz{Er ist die Trefferzahl, die bei jedem Durchgang herauskommt.} \\ \kreuz{Er ist der Durchschnitt der Trefferzahlen vieler Durchgänge.} \\ \kreuz{Er ist die Höhe der höchsten Säule.}

\binomialverteilung{10}{0.5}
\teil Das Diagramm zeigt die Verteilung einer binomialverteilten Zufallsgröße $X$ mit $n = 10$; der Erwartungswert von $X$ ist ganzzahlig. Lies ihn ab und bestimme $p$.

\binomialverteilung{10}{0.3}
$p = $ \leerfeld
\teil Die Diagramme zeigen die Verteilungen von $X$ und $Y$, beide binomialverteilt mit $n = 10$ und ganzzahligem Erwartungswert. Welche hat die größere Trefferwahrscheinlichkeit? Bestimme beide.

\binomialverteilung{10}{0.2} \binomialverteilung{10}{0.6}
$p_X = $ \leerfeld , $p_Y = $ \leerfeld
\end{teile}
\end{aufgabe}

\begin{aufgabe}{Diagramm – Prüfungsaufgaben}
\begin{teile}
\teil Das Diagramm zeigt die Verteilung einer binomialverteilten Zufallsgröße $X$ mit $n = 30$ und unbekanntem $p$; der Erwartungswert von $X$ ist ganzzahlig. Bestimme $p$ \hfill (Abitur 2026 GK).

\binomialverteilung{30}{0.2}
$p = $ \leerfeld
\teil $X$ ist binomialverteilt mit unbekanntem $n$ und der Trefferwahrscheinlichkeit 50\,\%. Das Diagramm zeigt zwei gleich hohe höchste Säulen bei 7 und 8. Begründe, dass $n$ nicht gerade ist \hfill (Abitur 2025 GK).

\binomialverteilung{15}{0.5}
\end{teile}
\end{aufgabe}
